\originalchapter{功、势能与能量}\label{ch:energy}
\originalsection{功与动能定理}
力沿位移方向推动物体会改变速率, 垂直速度的力改变方向却不瞬时改变速率. 这提示用点积表达力的速度作用. 对恒质量质点,
\begin{equation}
\vct F\cdot\vct v=m\dot{\vct v}\cdot\vct v
=\frac{\mathrm d}{\mathrm dt}\left(\tfrac12mv^2\right).
\end{equation}
功定义为 $W=\int_C\vct F\cdot\dd\vct r$, 路径 $C$ 是实际运动路径, 从上式积分得到 $\Delta K=W_{\rm net}$, $K=mv^2/2$. 恒力直线位移时 $W=\vct F\cdot\Delta\vct r$. 位置变力不能用某个末端力乘总距离. 功率 $P=\mathrm dW/\mathrm dt=\vct F\cdot\vct v$, 单位 W=J/s.

对两粒子, 总动能变化为各外力功和内力功之和. 虽然内力 $\vct F_{12}=-\vct F_{21}$, 位移一般不同, 所以内力总功为 $\vct F_{12}\cdot(\dd\vct r_2-\dd\vct r_1)$, 未必为零. 弹簧恢复力可以把内部储能转成动能. 不能把“内力不改变总动量”误写成“内力不做功”.

\originalsection{保守力与势能}
若力的功只取决于起终点, 称该力为保守力. 定义势能差 $U(B)-U(A)=-\int_A^B\vct F\cdot\dd\vct r$. 加上任意常数不改变力. 一维中 $F_x=-\mathrm dU/\mathrm dx$; 三维中 $\vct F=-\nabla U$, 即每个方向的力是该方向势能下降率. 在适当单连通区域内, 连续可微力场的旋度为零可用于检查保守性; 有奇点或非单连通区域时还需检查闭路积分.

地表重力势能取 $U=mgy$, 弹簧为 $U=kx^2/2$, 引力为 $U=-GMm/r$ 并约定无穷远为零. 若仅有这些保守力, 动能定理给 $\Delta(K+U)=0$. 若有非保守外力, $\Delta(K+U)=W_{\rm nc}$. 摩擦把宏观机械能转成内能, 所以机械能减少并不意味着总能量消失.

内部摩擦也可减少系统机械能, 如两个物体粘合碰撞. 把机械能、内能与燃料或人体的源能量一起记账, 才能说总能量守恒. 电机提升重物时输入电能包括重力势能增加、动能变化和耗散. 本册仅介绍所需的能量账目, 热力学状态、热与温度的完整理论另属热学册.

\originalsection{势能曲线与平衡}
给定总能量 $E$, 一维允许运动区满足 $K=E-U(x)\ge0$. $E=U$ 为转向点, 其间速率为 $\sqrt{2(E-U)/m}$. 平衡点要求 $U'(x_e)=0$. 稳定小振动还需 $U''(x_e)>0$, 因为泰勒展开给 $F\simeq-U''(x_e)(x-x_e)$, 位移产生恢复力. 若 $U''<0$, 扰动受到同向力而放大; 若 $U''=0$, 要看更高阶项, 不能立即判为中性平衡.

从能量关系可得旅行时间 $t=\int\sqrt{m/[2(E-U)]}\dd x$, 在单调运动段选正向时间. 转向点处速率为零并不意味着时间一定发散, 要看根号奇点是否可积. 能量方程给速率大小, 方向须由初值和转向行为确定.

\begin{example}
小块从固定光滑球顶静止滑下, 球半径 $R$. 在偏离顶点角 $\theta$ 处求速率和支持力.
\end{example}
\begin{solution}
高度下降 $R(1-\cos\theta)$, 所以 $v^2=2gR(1-\cos\theta)$. 向球心投影: $mg\cos\theta-N=mv^2/R$, 得 $N=mg(3\cos\theta-2)$. 失去接触在 $\cos\theta=2/3$. 此后的路径是抛物线, 不能继续把径向距离当成固定 $R$.
\end{solution}

\originalsection{非保守力的判断}
例如平面力 $\vct F=(ay)\uv x+(ax)\uv y$ 可取 $U=-axy$, 所以虽各分量同时依赖两个坐标, 它仍保守. 若改为 $\vct F=(ay)\uv x-(ax)\uv y$, 沿边长 $l$ 的方形闭路积分一般非零, 从 $\partial F_y/\partial x=-a$ 与 $\partial F_x/\partial y=a$ 也可看出不保守. 保守性的判断针对力场及定义域, 不是针对某一次路径上恰好零功.
\originalsection{路径功的坐标计算}
在直角坐标中, $W=\int(F_x\dd x+F_y\dd y+F_z\dd z)$. 若路径以参数 $s$ 表示, 则把 $x(s),y(s),z(s)$ 及其导数全部代入同一积分. 在柱坐标中, 微小位移为 $\dd\vct r=\uv r\dd r+\uv\theta r\dd\theta+\uv z\dd z$, 所以
\begin{equation}
W=\int(F_r\dd r+F_\theta r\dd\theta+F_z\dd z).
\end{equation}
$\dd\theta$ 本身不是弧长, 必须乘 $r$. 例如半径 $R$ 的圆弧上恒切向力 $F_\theta=F_0$ 从角 0 到 $\Theta$, 做功 $F_0R\Theta$; 径向力在同弧上不做功.

三维势能的偏导 $\partial U/\partial x$ 表示保持 $y,z$ 固定时沿 $x$ 的变化率, 梯度 $\nabla U$ 收集三个偏导. 旋度零判据在本册实际用于分量相等检查, 如平面条件 $\partial F_x/\partial y=\partial F_y/\partial x$. 完整向量微积分可接电磁学册; 本册已给判断所需的含义与域条件.
